\subsection{Limitations}
\label{sec:future-limit}

\paragraph{Network conditions.}
Mode A (bridge-coalesced) and Mode B (paced \SI{100}{\micro\second})
bound the realistic spectrum but do not cover every deployment.
Particularly: cloud LB-to-pod paths sometimes apply per-flow
shapers that approximate Mode B even without an explicit attacker
pacing strategy. We have not measured those configurations
directly.

\paragraph{Single-CPU constraint.}
Every server runs in a \texttt{--cpus=1} container. This is the
single-CPU worst case. Multi-CPU servers fan out work across cores,
but the per-connection cost is still serial per the affected
worker, so multi-CPU does not change the per-connection ratio --
it changes how many concurrent attackers are needed to saturate
the deployment.

\paragraph{TLS overhead not measured.}
Every test is plain HTTP. TLS adds AEAD framing on top of TCP;
TLS records buffer multiple bytes by their nature, which could
incidentally aggregate chunked-TE chunks at the record boundary.
We have not measured the magnitude.

\paragraph{HTTP/2 and HTTP/3 not measured (Wave 1+2).}
The 27 servers in our sample were measured under HTTP/1.1
chunked-TE only. HTTP/2 uses DATA frames with explicit length
fields; HTTP/3 uses QUIC streams with similar framing. Whether
the per-chunk amplification class transfers, mitigates, or
manifests differently under those protocols is unknown.

\paragraph{Cold start vs warm path.}
JVM tests in particular include some JIT warm-up cost in the
first few requests. We report the warm-path numbers (request
2--N of each cell); cold-path measurements would change the
absolute timings but not the asymptotic shape.

\paragraph{Production traffic distribution.}
The probability that real-world traffic includes per-chunk
amplification attempts is unmeasured. We have not deployed
honeypots or instrumented production servers.

\subsection{Future work}
\label{sec:future-work}

\paragraph{Wave 3: HTTP/2, HTTP/3, WebSocket, gRPC.}
We plan to measure the same 27 servers under HTTP/2 multiplexed
DATA frames, HTTP/3 over QUIC streams, WebSocket per-frame
amplification, and gRPC streaming. Preliminary expectation: the
class transfers to all four, with HTTP/2 \emph{worsening} the
per-connection cost via stream multiplexing and HTTP/3 partially
mitigating via QUIC's flow control.

\paragraph{Multipart per-part amplification.}
\fwid{multipart/form-data} is structurally similar: each part has
a header section delivered to the parser callbacks per chunk.
Companion work in
\texttt{findings/starlette/2026-05/STARLETTE\_FORMPARSER\_QUADRATIC/}
(retracted as a security finding after patch verification but the
underlying parser shape is the same) suggests the class transfers.

\paragraph{Production telemetry.}
We have no data on what fraction of production HTTP traffic
includes chunked-TE bodies with small chunks. A honeypot study or
a partnership with a high-traffic CDN would let us estimate the
real-world prevalence of the attack pattern.

\paragraph{Frontend buffering measurement at CDNs.}
Cloudflare, Fastly, Akamai, AWS ALB, and GCP HTTP(S) LB each have
proprietary chunked-body handling. None publicly document whether
they aggregate before forwarding to origin. A measurement campaign
across these would significantly improve the operational
recommendation in Section~\ref{sec:mit-summary}.

\paragraph{Formal model of batching correctness.}
We have argued informally (Section~\ref{sec:rc-spec}) that
batching is spec-compliant under RFC 9112 and ASGI / WSGI /
Servlet. A formal model -- expressed in TLA+ or a similar
specification language -- would make the argument bulletproof
and would let server maintainers verify their own batching
implementations.

\paragraph{Detection signatures.}
Building an open-source IDS / WAF signature for the attack
pattern is straightforward (Section~\ref{sec:tm-detect}) but
requires balancing false positives against attacker evasion.
A reference Snort / Suricata rule set, validated on our
reproducer harness, would be a useful publishable artifact.
